\documentclass[11pt,a4paper]{article}
\usepackage[utf8]{inputenc}
\usepackage[T1]{fontenc}
\usepackage{amsmath,amssymb,bm,mathtools}
\usepackage{booktabs,multirow,array,tabularx,longtable}
\usepackage{graphicx}
\usepackage[colorlinks=true,citecolor=blue,linkcolor=blue,urlcolor=blue]{hyperref}
\usepackage{geometry}
\usepackage[numbers,sort&compress]{natbib}
\usepackage{xcolor}
\usepackage{caption}
\usepackage{subcaption}
\usepackage{enumitem}
\usepackage{setspace}
\usepackage{parskip}
\usepackage{colortbl}
\geometry{top=1in,bottom=1in,left=1.1in,right=1.1in}

\newcolumntype{C}[1]{>{\centering\arraybackslash}p{#1}}
\newcolumntype{L}[1]{>{\raggedright\arraybackslash}p{#1}}

\definecolor{ours}{HTML}{d4edda}
\definecolor{header}{HTML}{f0f0f0}
\definecolor{sota}{HTML}{ffeeba}

\title{\LARGE\bfseries EB-JEPA: Energy-Based Joint Embedding Predictive Architecture\\
for Self-Supervised EEG Representation Learning\\[4pt]
{\large A Technical Report on the TUAB Abnormality Benchmark}}

\author{
  Yani Hammache$^{\dagger}$ \quad
  Th\'eo Basseras$^{\dagger}$ \quad
  Titouan Vasnier$^{\dagger}$ \quad
  Timothy Courtois$^{\dagger}$\\[4pt]
  \small$^{\dagger}$Hackathon Team, Universit\'e Paris-Saclay / Dalia Cluster
}
\date{June 2026}

\begin{document}
\maketitle

%% ============================================================
\begin{abstract}
%% ============================================================
We study self-supervised representation learning for clinical EEG through an
Energy-Based Joint Embedding Predictive Architecture (EB-JEPA).  Two independently
corrupted views of the same 10-second EEG window are aligned in latent space by an
energy function combining invariance, variance, and covariance regularisation
(VICReg~\cite{vicreg}) or a stronger Gaussianisation surrogate
(SIGReg / Batched Characteristic Slicing~\cite{epps1983}).
Four encoder families are evaluated on the TUH Abnormal EEG Corpus
(TUAB v3.0.1~\cite{tuab}, 2717 training and 276 evaluation recordings,
binary normal/abnormal classification):
a temporal Conv1D network (0.4M parameters), a patch Transformer in the style of
LaBraM~\cite{labram} (1.2M), an FFT-tokenised Transformer in the style of
BIOT~\cite{biot} (1.2M), and a hierarchical channel-pool Transformer in the
style of EEGPT~\cite{eegpt} (0.8M).

Our best frozen encoder (Conv1D + corruption augmentation + SIGReg) reaches
\textbf{0.775 balanced accuracy (BACC) per-window} and \textbf{0.825 BACC per-recording}
(mean-pooling 16 window embeddings per patient), matching ContraWR and approaching
LaBraM-Base (0.814 per-window).  Fine-tuning adds 1.2 pp to reach \textbf{0.837 BACC
per-recording} and AUROC \textbf{0.919}.

The central contribution of this report is \emph{diagnostic}:
we identify and document three systematic failure modes ---
a VICReg coefficient bug ($\lambda{=}1$ vs.\ the paper-correct $\lambda{=}25$),
an architectural conflict between EEGPT's channel-pooling and channel-drop
augmentation, and a per-window vs.\ per-recording evaluation gap
(${\approx}5$\,pp for the same encoder) that initially caused an inflated
performance claim.  We describe how each failure was diagnosed, quantified, and
corrected.  An auxiliary spectral regularisation loss further yields the strongest
per-recording frozen result (0.836 BACC), leaving open whether it provides
independent signal or compensates for the underweighted invariance term.
\end{abstract}

\tableofcontents
\newpage

%% ============================================================
\section*{Introduction}
\addcontentsline{toc}{section}{Introduction}
\label{sec:intro}
%% ============================================================

This report documents the application of Energy-Based Joint Embedding Predictive
Architectures (EB-JEPA,~\cite{ebjepa_lib}) to clinical EEG abnormality detection on
the TUH Abnormal EEG Corpus (TUAB).  The work was conducted during a 48-hour
hackathon at the Dalia cluster, under the broader EB-JEPA research initiative.

\paragraph{Structure of the report.}
\begin{description}[noitemsep]
  \item[Section~\ref{sec:sota}: State of the Art.]
    We first clarify the three distinct \emph{scoring styles} used across the EEG
    literature (per-window, per-recording via window aggregation, per-recording from
    full recording), as confusing them produces spurious comparison advantages.
    We then provide the complete \emph{provenance table} of published results and
    our own ground-truth numbers from the repository, covering supervised ConvNets,
    masked SSL (BIOT, LaBraM), contrastive SSL (BENDR, ContraWR), and recent
    foundation models (FEMBA~\cite{femba}, EEG-Bench~\cite{eeg_bench}).

  \item[Section~\ref{sec:jepa}: EB-JEPA Results.]
    We present the two-view energy concept underlying EB-JEPA, describe the TUAB
    dataset and the critical distinction between SSL training and probe evaluation
    splits, and then report all experimental results from the hackathon in a single
    comprehensive table.

  \item[Section~\ref{sec:math}: Per-Model Mathematical Descriptions.]
    For every encoder family tested, we provide the complete energy function with
    explicit loss weights, encoder architecture details, and the MLP probe
    specification.  We then describe the training and evaluation protocol exactly
    as implemented.

  \item[Section~\ref{sec:best}: Combined Best Results.]
    We consolidate the best numbers from the literature and from this work,
    filtered to per-recording metrics only (the clinically meaningful unit of
    analysis for TUAB).

  \item[Section~\ref{sec:conclusion}: Conclusion and Future Work.]
    We conclude with open scientific questions, methodological lessons, and
    a prioritised list of future experiments.
\end{description}

\paragraph{Research question.}
Can a small (0.4M parameter) encoder pre-trained \emph{without labels} on 2717
unlabelled EEG recordings learn representations that generalise to held-out
patients and match or exceed supervised models trained with labels?

%% ============================================================
\section{State of the Art on TUAB EEG Abnormality Detection}
\label{sec:sota}
%% ============================================================

\subsection{Scoring Styles}
\label{sec:scoring}

Published results on TUAB use heterogeneous evaluation conventions.
We precisely define four \emph{scoring styles} used in the literature:

\paragraph{(S1) Per-window (WIN).}
The evaluation set is segmented into all non-overlapping 10-second windows.
Each window is independently classified; it inherits the label of its parent
recording.  BACC is computed over all individual window predictions.
For TUAB eval (276 recordings, 30--120 min each), this yields ${\approx}400$k
test samples.
This is the convention used by BIOT~\cite{biot}, LaBraM~\cite{labram},
ContraWR~\cite{contrawr}, and most published SSL models.  The test samples
are \emph{not independent}: windows from the same recording share the same
label and much of the same signal, inflating effective sample size.

\paragraph{(S2) Per-recording via window embedding aggregation (REC-WIN).}
A fixed number $N$ of evenly-spaced windows are extracted per recording.
Each window is encoded separately; the $N$ embeddings are \emph{mean-pooled}
into one vector.  The probe classifies this pooled vector; one prediction
is made per recording.  This is the default mode of our framework with $N{=}16$.
The noise-averaging effect of pooling systematically raises BACC by ${\approx}5$\,pp
over WIN for the same encoder and probe.

\paragraph{(S3) Per-recording via majority vote (REC-VOTE).}
Similar to REC-WIN but aggregates \emph{decisions} rather than embeddings:
each window is independently classified, and the recording label is the
majority class.  This style is used in some clinical pipelines but not in the
SSL literature surveyed here.

\paragraph{(S4) Per-recording from the full recording (REC-FULL).}
The entire recording (variable length, up to 120 min) is fed as a single
input to a model capable of handling variable-length sequences (e.g., an SSM
or a hierarchical model).  No windowing or aggregation is required.
None of our models use this mode; it is listed for completeness.

\medskip
\noindent\textbf{Critical warning.}
Comparing a REC-WIN or REC-FULL number directly to a WIN number as if they
were the same quantity produces a spurious ${\approx}5$\,pp advantage to the
per-recording model.  Section~\ref{sec:conclusion} describes how this gap was
inadvertently exploited and then corrected in this project.

\subsection{State-of-the-Art Provenance Table}
\label{sec:sota_table}

Table~\ref{tab:sota} compiles every TUAB result we could verify against a
primary source.  We call this the \emph{provenance table}: every entry is
annotated with its source paper, so the reader can verify the number directly.
Entries in~\colorbox{sota}{yellow} denote results from the EEG-Bench
benchmark~\cite{eeg_bench} or FEMBA~\cite{femba} added in this update.

\begin{table}[htbp]
\centering
\caption{%
  Provenance table: state-of-the-art results on TUAB EEG abnormality detection.
  WIN = per-window; REC-WIN = per-recording via window embedding pooling;
  REC-FT = per-recording, fine-tuned end-to-end.
  ``Frozen'' = frozen encoder + linear/MLP probe.
  ``FT'' = fine-tuned encoder.
  $^{a}$Accuracy (not BACC), early eval script.
  $^{b}$Reported as accuracy; TUAB eval split is near-balanced
  (54\%\,normal, 46\%\,abnormal), so acc $\approx$ BACC within 2\,pp.
  $^{c}$AUROC estimated from ROC curve in the paper.
  $^{d}$Protocol unclear; likely REC-WIN with averaging over non-overlapping chunks
  before classification.
  $^{\star}$AUROC measured under REC-WIN ($N{=}16$); WIN AUROC not recorded.
  VICReg$^{*}$\,=\,buggy \texttt{inv\_coeff=1}; VICReg\,=\,corrected \texttt{inv\_coeff=25}.
}
\label{tab:sota}
\scriptsize
\setlength{\tabcolsep}{3pt}
\resizebox{\textwidth}{!}{%
\begin{tabular}{L{2.0cm}L{3.2cm}C{1.1cm}C{1.1cm}C{1.1cm}L{2.6cm}L{1.3cm}}
\toprule
\textbf{Scoring} &
\textbf{Model} &
\textbf{BACC} &
\textbf{AUROC} &
\textbf{Params} &
\textbf{Source} &
\textbf{Mode} \\
\midrule
\multicolumn{7}{l}{\textit{Floor}} \\
WIN & Random encoder & 0.52 & -- & -- & This work & Frozen \\
\midrule
\multicolumn{7}{l}{\textit{Supervised baselines}} \\
WIN & ShallowConvNet~\cite{shallowconvnet} & 0.777 & 0.857 & 44k & Schirrmeister 2017 & FT \\
WIN & EEGNet~\cite{eegnet} & 0.796 & 0.882 & 9k & Kiessner 2023~\cite{kiessner2023} & FT \\
WIN & Deep4Net~\cite{shallowconvnet} & 0.816$^{b}$ & -- & -- & Schirrmeister 2018 & FT \\
\midrule
\multicolumn{7}{l}{\textit{Self-supervised foundation models (published)}} \\
WIN & ContraWR~\cite{contrawr} & 0.775 & -- & -- & Yang 2021 & Frozen \\
WIN & BENDR~\cite{bendr} & 0.717 & -- & 4M & Kastrati 2025~\cite{eeg_bench} & FT$^{d}$ \\
WIN & BIOT~\cite{biot} & 0.802 & -- & 3.2M & Yang 2023 & Frozen \\
WIN & BIOT~\cite{biot} & 0.802$^{c}$ & -- & 3.2M & Bindra 2026~\cite{spectral_audit} & FT$^{d}$ \\
WIN & LaBraM-Base~\cite{labram} & 0.814 & 0.890$^{c}$ & 5.8M & Jiang 2024 & Frozen \\
WIN & LaBraM-Base~\cite{labram} & 0.786 & -- & 5.8M & Bindra 2026~\cite{spectral_audit} & FT$^{d}$ \\
\cellcolor{sota}REC-FT & \cellcolor{sota}LaBraM-Base~\cite{labram} & \cellcolor{sota}0.838 & \cellcolor{sota}-- & \cellcolor{sota}5.8M & \cellcolor{sota}Kastrati 2025~\cite{eeg_bench} & \cellcolor{sota}FT \\
\cellcolor{sota}REC-FT & \cellcolor{sota}Neuro-GPT~\cite{eeg_bench} & \cellcolor{sota}0.696 & \cellcolor{sota}-- & \cellcolor{sota}-- & \cellcolor{sota}Kastrati 2025~\cite{eeg_bench} & \cellcolor{sota}FT \\
WIN & EEGPT~\cite{eegpt} & 0.801 & -- & 10M & Bindra 2026~\cite{spectral_audit} & FT$^{d}$ \\
\cellcolor{sota}REC-FT & \cellcolor{sota}FEMBA-Base~\cite{femba} & \cellcolor{sota}0.8105 & \cellcolor{sota}0.8894 & \cellcolor{sota}47.7M & \cellcolor{sota}Tegon 2025 & \cellcolor{sota}FT \\
\cellcolor{sota}REC-FT & \cellcolor{sota}FEMBA-Large~\cite{femba} & \cellcolor{sota}0.8147 & \cellcolor{sota}0.8892 & \cellcolor{sota}77.8M & \cellcolor{sota}Tegon 2025 & \cellcolor{sota}FT \\
\cellcolor{sota}REC-FT & \cellcolor{sota}FEMBA-Huge~\cite{femba} & \cellcolor{sota}\textbf{0.8182} & \cellcolor{sota}\textbf{0.9005} & \cellcolor{sota}386M & \cellcolor{sota}Tegon 2025 & \cellcolor{sota}FT \\
\midrule
\multicolumn{7}{l}{\textit{\colorbox{ours}{EB-JEPA (this work)} -- per-window (WIN), frozen encoder}} \\
WIN & \cellcolor{ours}Conv1D + VICReg$^{*}$ & \cellcolor{ours}0.756 & \cellcolor{ours}0.888$^{\star}$ & \cellcolor{ours}0.4M & \cellcolor{ours}This work & \cellcolor{ours}Frozen \\
WIN & \cellcolor{ours}Conv1D + VICReg$^{*}$ + corrupt & \cellcolor{ours}0.770 & \cellcolor{ours}0.904$^{\star}$ & \cellcolor{ours}0.4M & \cellcolor{ours}This work & \cellcolor{ours}Frozen \\
WIN & \cellcolor{ours}Conv1D + SIGReg + corrupt & \cellcolor{ours}0.775 & \cellcolor{ours}0.904$^{\star}$ & \cellcolor{ours}0.4M & \cellcolor{ours}This work & \cellcolor{ours}Frozen \\
WIN & \cellcolor{ours}Conv1D + VICReg$^{*}$ + spectral 0.1 & \cellcolor{ours}0.765 & \cellcolor{ours}0.887$^{\star}$ & \cellcolor{ours}0.4M & \cellcolor{ours}This work & \cellcolor{ours}Frozen \\
WIN & \cellcolor{ours}LaBraM-style + VICReg$^{*}$ & \cellcolor{ours}0.775$^{a}$ & \cellcolor{ours}-- & \cellcolor{ours}1.2M & \cellcolor{ours}This work & \cellcolor{ours}Frozen \\
WIN & \cellcolor{ours}BIOT-style + VICReg$^{*}$ & \cellcolor{ours}0.775$^{a}$ & \cellcolor{ours}-- & \cellcolor{ours}1.2M & \cellcolor{ours}This work & \cellcolor{ours}Frozen \\
WIN & \cellcolor{ours}BIOT-style + SIGReg & \cellcolor{ours}0.783$^{a}$ & \cellcolor{ours}-- & \cellcolor{ours}1.2M & \cellcolor{ours}This work & \cellcolor{ours}Frozen \\
WIN & \cellcolor{ours}EEGPT-style + VICReg$^{*}$ & \cellcolor{ours}collapsed & \cellcolor{ours}-- & \cellcolor{ours}0.8M & \cellcolor{ours}This work & \cellcolor{ours}Frozen \\
\midrule
\multicolumn{7}{l}{\textit{\colorbox{ours}{EB-JEPA (this work)} -- per-recording via window pooling (REC-WIN, $N{=}16$)}} \\
REC-WIN & \cellcolor{ours}Conv1D + VICReg$^{*}$ & \cellcolor{ours}0.796 & \cellcolor{ours}0.888 & \cellcolor{ours}0.4M & \cellcolor{ours}This work & \cellcolor{ours}Frozen \\
REC-WIN & \cellcolor{ours}Conv1D + VICReg$^{*}$ + corrupt & \cellcolor{ours}0.825 & \cellcolor{ours}0.904 & \cellcolor{ours}0.4M & \cellcolor{ours}This work & \cellcolor{ours}Frozen \\
REC-WIN & \cellcolor{ours}Conv1D + SIGReg + corrupt & \cellcolor{ours}0.825 & \cellcolor{ours}0.904 & \cellcolor{ours}0.4M & \cellcolor{ours}This work & \cellcolor{ours}Frozen \\
REC-WIN & \cellcolor{ours}Conv1D + VICReg$^{*}$ + spectral 0.1 & \cellcolor{ours}\textbf{0.836} & \cellcolor{ours}0.887 & \cellcolor{ours}0.4M & \cellcolor{ours}This work & \cellcolor{ours}Frozen \\
REC-WIN & \cellcolor{ours}Conv1D + VICReg$^{*}$ + corrupt (FT) & \cellcolor{ours}0.812 & \cellcolor{ours}\textbf{0.908} & \cellcolor{ours}0.4M & \cellcolor{ours}This work & \cellcolor{ours}FT \\
\bottomrule
\end{tabular}}
\end{table}

\paragraph{Key observations from Table~\ref{tab:sota}.}
\begin{enumerate}[noitemsep]
  \item The best published TUAB BACC is held by FEMBA-Huge at 81.82\% (fine-tuned,
    trained on 21,000h of unlabelled EEG from TUEG), followed by LaBraM-Large/Huge
    (82.26--82.58\% in FEMBA's evaluation; 83.8\% fine-tuned in EEG-Bench).
  \item Among methods using the same TUAB training set without external pretraining,
    LaBraM-Base (0.814 frozen, WIN) is the primary target for our approach.
  \item Our best frozen WIN BACC (0.775, SIGReg+corrupt) matches ContraWR and falls
    3.9\,pp below LaBraM-Base, with only 0.4M parameters and 20 training epochs.
  \item The spectral audit of Bindra \& Panwar~\cite{spectral_audit} reveals that
    supervised and SSL models on TUAB rely substantially on the aperiodic (1/f)
    broadband spectral slope ($\Delta$BACC 0.07--0.13 after flattening), which may
    partially confound abnormality detection with age-related spectral changes.
    EB-JEPA's SIGReg Gaussianisation surrogate is architecturally designed to
    prevent representational collapse, but whether it avoids aperiodic shortcuts
    remains an open question.
\end{enumerate}

%% ============================================================
\section{EB-JEPA Results}
\label{sec:jepa}
%% ============================================================

\subsection{EB-JEPA Two-View Energy Concept}
\label{sec:jepa_concept}

EB-JEPA is our EEG adaptation of the Joint Embedding Predictive Architecture
paradigm~\cite{ijepa,lecun2022}.  We adopt a \emph{two-view energy-based} formulation
adapted to time-series signals, which differs from the standard image JEPA in that
both views are corrupted versions of the \emph{same} window (rather than masked
spatial patches), and the encoder is trained to produce the same latent representation
for both views.

\paragraph{Two-view concept with two corruption pathways.}
Given a 10-second EEG window $\mathbf{x} \in \mathbb{R}^{C \times T}$
($C{=}19$ channels, $T{=}2000$ samples at 200\,Hz), two independently corrupted
views are produced as follows:
\begin{itemize}[noitemsep]
  \item \textbf{Corruption pathway (a):} Apply augmentation stack $\mathcal{A}_a$ to
    $\mathbf{x}$ to obtain $\tilde{\mathbf{x}}_a = \mathcal{A}_a(\mathbf{x})$.
  \item \textbf{Corruption pathway (b):} Apply independently sampled augmentation
    stack $\mathcal{A}_b$ to obtain $\tilde{\mathbf{x}}_b = \mathcal{A}_b(\mathbf{x})$.
\end{itemize}
Both $\mathcal{A}_a$ and $\mathcal{A}_b$ are drawn from the same family but with
independently sampled random parameters (noise seed, masked channels, etc.).
A shared encoder $f_\theta$ maps each view to a latent vector:
\begin{equation}
  \mathbf{z}_a = g_\phi(f_\theta(\tilde{\mathbf{x}}_a)), \quad
  \mathbf{z}_b = g_\phi(f_\theta(\tilde{\mathbf{x}}_b)), \quad
  \mathbf{z}_a, \mathbf{z}_b \in \mathbb{R}^{d_p},\; d_p{=}1024,
\end{equation}
where $g_\phi$ is a projector MLP ($d \to 1024 \to 1024$) that maps from the
encoder output dimension $d{=}256$ to the projector space.
The energy function $E(\mathbf{z}_a, \mathbf{z}_b) = \|\mathbf{z}_a - \mathbf{z}_b\|^2$
penalises disagreement between the two views.
Without an anti-collapse mechanism, the trivial global minimum
$\mathbf{z}_a = \mathbf{z}_b = \mathbf{0}$ satisfies $E{=}0$ with zero information.
VICReg~\cite{vicreg} and SIGReg prevent this collapse by adding variance and
covariance regularisation terms (described formally in Section~\ref{sec:math_ssl}).

The two-view concept is illustrated in Figure~\ref{fig:ebjepa}: the world model
(encoder) is trained to map two different corruptions of the same EEG brain state to
the same point in latent space.

\begin{figure}[htbp]
\centering
\fbox{\parbox{0.90\textwidth}{\centering\vspace{8pt}
$\mathbf{x} \in \mathbb{R}^{19 \times 2000}$ (EEG window, 10s)
\\[6pt]
$\downarrow \quad \mathcal{A}_a(\cdot) \qquad\qquad\qquad \mathcal{A}_b(\cdot) \quad \downarrow$
\\[4pt]
$\tilde{\mathbf{x}}_a \xrightarrow{f_\theta \circ g_\phi} \mathbf{z}_a \in \mathbb{R}^{1024}
\qquad \mathbf{z}_b \in \mathbb{R}^{1024} \xleftarrow{f_\theta \circ g_\phi} \tilde{\mathbf{x}}_b$
\\[6pt]
$E = \underbrace{\|\mathbf{z}_a - \mathbf{z}_b\|^2}_{\text{invariance (energy)}}
+ w_V \underbrace{L_V(\mathbf{z}_a, \mathbf{z}_b)}_{\text{variance (anti-collapse)}}
+ w_C \underbrace{L_C(\mathbf{z}_a, \mathbf{z}_b)}_{\text{covariance (decorrelation)}}$
\vspace{8pt}
}}
\caption{EB-JEPA two-view training scheme.  Two independently corrupted views of the
same EEG window (pathways a and b) are mapped to the same latent space through a shared
encoder $f_\theta$ and projector $g_\phi$.  The energy function combines invariance
with variance and covariance regularisation to prevent representational collapse.
The encoder acts as the \emph{world model}: it learns to abstract away the specific
corruption and retain the clinically relevant brain state.}
\label{fig:ebjepa}
\end{figure}

\subsection{EEG Dataset: SSL Training vs.\ Probe Evaluation Splits}
\label{sec:dataset}

All experiments use the \textbf{TUH Abnormal EEG Corpus (TUAB) v3.0.1}~\cite{tuab}.
A critical distinction exists between the \emph{SSL split} (used during pre-training)
and the \emph{probe split} (used to train and evaluate the downstream classifier).

\paragraph{Recording configuration.}
19-channel scalp EEG in the international 10-20 system (FP1, FP2, F3, F4, C3, C4,
P3, P4, O1, O2, F7, F8, T3, T4, T5, T6, FZ, CZ, PZ), in a fixed channel order
across every file.  Recordings span 10--120 minutes (clinical monitoring sessions).
The binary label (\emph{normal} vs.\ \emph{abnormal}) is assigned by a board-certified
neurologist to the entire recording.

\paragraph{Offline preprocessing (applied once, provided by the dataset).}
The TUAB\_PREPROCESSED folder contains signals already processed by the dataset
providers in the following order:
\begin{enumerate}[noitemsep]
  \item \textbf{Channel selection.}  19 scalp electrodes in the 10-20 system,
    fixed order across all files.  Raw TUAB recordings vary in the number of
    channels; only the 19 standard channels are retained.
  \item \textbf{Resampling to 200\,Hz.}  Raw recording rates vary (250, 256,
    512\,Hz depending on the acquisition system).  All signals are resampled to a
    common 200\,Hz grid.  At 200\,Hz a 10-second window contains exactly 2000
    samples per channel, which divides cleanly for strided convolutions and patch
    encoders.
  \item \textbf{60\,Hz notch filter.}  TUAB is a US dataset; US mains frequency
    is 60\,Hz (not 50\,Hz as in Europe).  A narrow notch removes this line-noise
    artefact before any learning.
  \item \textbf{0.1--75\,Hz bandpass.}  The standard clinical EEG band, covering
    delta (0.5--4\,Hz), theta (4--8\,Hz), alpha (8--13\,Hz), beta (13--30\,Hz),
    and low gamma (30--70\,Hz) while discarding DC drift ($<$0.1\,Hz) and
    high-frequency noise ($>$75\,Hz, comfortably below the 100\,Hz Nyquist at
    200\,Hz).  No normalization is applied at this stage; amplitudes remain in
    native \textmu{}V.
\end{enumerate}

\paragraph{Online preprocessing (applied at load time, per window).}
After the offline pipeline, each 10-second window loaded during training or
evaluation is \textbf{z-scored per channel}:
\[
  \hat{x}_c = \frac{x_c - \mu_c}{\sigma_c + 10^{-6}},
\]
where $\mu_c$ and $\sigma_c$ are the mean and standard deviation of channel $c$
computed \emph{within the window only} (no statistics leak across windows or
patients).  This removes per-channel amplitude offsets caused by electrode
impedance variability and brings all channels to a common scale regardless of the
underlying µV amplitude.  No global normalization (across recordings or the full
dataset) is used.

\paragraph{Patient-disjoint splits.}
TUAB provides an official patient-disjoint train/eval split:
\begin{itemize}[noitemsep]
  \item \textbf{SSL training split} (2\,717 recordings: 1\,385 normal, 1\,332 abnormal):
    used for SSL pre-training \emph{without labels}.  The encoder $f_\theta$ sees only
    the raw EEG signal; the normal/abnormal labels are discarded.
  \item \textbf{Evaluation split} (276 recordings: 150 normal, 126 abnormal):
    used exclusively for final performance estimation.  No hyperparameter selection is
    performed on this split.
\end{itemize}
No recording from a given patient appears in both splits, preventing
patient-level data leakage.

\paragraph{SSL pre-training (ssl split).}
The encoder is pre-trained using the SSL training split.  Labels are discarded.
The two-view SSL objective (VICReg or SIGReg) is optimised over 20 epochs on
non-overlapping 10-second windows extracted from all 2\,717 training recordings.

\paragraph{Probe training (probe split = ssl training split with labels).}
After SSL pre-training, the encoder is frozen.  An MLP probe is trained on the
\emph{same} 2\,717 training recordings, now using the normal/abnormal labels.
For each recording, $N{=}16$ evenly-spaced 10-second windows are extracted and
encoded; their embeddings are mean-pooled into one 256-dim vector (REC-WIN protocol).
The probe is trained on these pooled vectors.

\paragraph{Comparison with the published SSL protocol.}
The published LaBraM, BIOT, and ContraWR evaluations use the per-window (WIN)
protocol on a larger eval split (${\approx}2339$ recordings, ${\approx}409$k test
windows).  Our eval split (276 recordings) is the official TUAB v3.0.1 split.
The different eval set sizes contribute to difficulty in direct comparison, beyond
the WIN vs.\ REC-WIN protocol gap.

\subsection{Full Results Table for JEPA Models Tested}
\label{sec:full_results}

Table~\ref{tab:full_results} reports all encoder $\times$ loss $\times$
augmentation combinations evaluated during the hackathon.

\begin{table}[htbp]
\centering
\caption{%
  Complete EB-JEPA results on TUAB v3.0.1 (276-recording eval split) and TUEV v2.0.0 (6-class event transfer).
  ``Frozen'' = frozen encoder, MLP probe.
  ``FT'' = entire network fine-tuned during probe training.
  WIN = per-window BACC on TUAB; REC = per-recording BACC on TUAB ($N{=}16$ windows, mean-pooled).
  BACC\textsubscript{TUEV} = 6-class macro balanced accuracy on TUEV event detection (frozen logistic probe, event-centred windows).
  ``collapsed'' = BACC $\leq$ random init floor.
  $^{a}$Accuracy (not BACC), early eval script.
  $^{\star}$AUROC under REC-WIN; WIN AUROC not recorded.
  VICReg$^{*}$ = buggy \texttt{inv\_coeff=1}; VICReg = corrected \texttt{inv\_coeff=25}.
  FT rows are 3-seed means at the \emph{final} epoch (no best-epoch selection).
}
\label{tab:full_results}
\small
\setlength{\tabcolsep}{4pt}
\resizebox{\textwidth}{!}{%
\begin{tabular}{L{4.0cm}C{1.3cm}C{1.3cm}C{1.3cm}C{1.3cm}C{1.3cm}L{1.5cm}L{1.4cm}C{1.2cm}}
\toprule
\textbf{Model / Config} &
\textbf{BACC\textsubscript{WIN}} &
\textbf{BACC\textsubscript{REC}} &
\textbf{AUROC\textsubscript{REC}} &
\textbf{BACC\textsubscript{TUEV}} &
\textbf{Score style} &
\textbf{Probe} &
\textbf{SSL loss} &
\textbf{Data} \\
\midrule
\multicolumn{9}{l}{\textit{Conv1D encoder (0.4M parameters) -- spectral regularisation ablation}} \\
Conv + VICReg$^{*}$ (base)          & 0.756 & 0.796 & 0.888 & 0.381 & WIN / REC-WIN & Frozen & VICReg$^{*}$ & SSL \\
Conv + VICReg$^{*}$ + spec 0.05     & --    & 0.821 & 0.884 & -- & REC-WIN & Frozen & VICReg$^{*}$ & SSL \\
Conv + VICReg$^{*}$ + spec 0.1      & 0.765 & 0.836 & 0.887 & \textbf{0.425} & WIN / REC-WIN & Frozen & VICReg$^{*}$ & SSL \\
Conv + VICReg$^{*}$ + spec 0.3      & --    & 0.789 & 0.877 & -- & REC-WIN & Frozen & VICReg$^{*}$ & SSL \\
Conv + VICReg$^{*}$ + spec 1.0      & --    & 0.785 & 0.874 & -- & REC-WIN & Frozen & VICReg$^{*}$ & SSL \\
Conv + VICReg$^{*}$ + fftc 0.05     & --    & 0.807 & 0.887 & -- & REC-WIN & Frozen & VICReg$^{*}$ & SSL \\
Conv + VICReg$^{*}$ + fftc 0.1      & --    & 0.798 & 0.881 & -- & REC-WIN & Frozen & VICReg$^{*}$ & SSL \\
Conv + VICReg$^{*}$ + fftc 0.3      & --    & 0.802 & 0.889 & -- & REC-WIN & Frozen & VICReg$^{*}$ & SSL \\
\midrule
\multicolumn{9}{l}{\textit{Conv1D encoder -- corruption augmentation ablation (3 seeds)}} \\
Conv + VICReg$^{*}$ + corrupt (s0)  & 0.770 & 0.825 & 0.904 & 0.382 & WIN / REC-WIN & Frozen & VICReg$^{*}$ & SSL \\
Conv + VICReg$^{*}$ + corrupt (s1)  & --    & 0.816 & 0.891 & -- & REC-WIN & Frozen & VICReg$^{*}$ & SSL \\
Conv + VICReg$^{*}$ + corrupt (s2)  & --    & 0.818 & 0.906 & -- & REC-WIN & Frozen & VICReg$^{*}$ & SSL \\
Conv + VICReg$^{*}$ + corrupt (big) & --    & 0.805 & 0.892 & -- & REC-WIN & Frozen & VICReg$^{*}$ & SSL \\
Conv + corrupt + spec 0.1           & --    & 0.802 & 0.899 & -- & REC-WIN & Frozen & VICReg$^{*}$ & SSL \\
Conv + SIGReg + corrupt             & 0.775 & 0.825 & 0.904 & 0.363 & WIN / REC-WIN & Frozen & SIGReg & SSL \\
\midrule
\multicolumn{9}{l}{\textit{Conv1D encoder -- fine-tuning, multi-corpus, and new objectives}} \\
Conv + VICReg$^{*}$ + corrupt (FT)  & --    & 0.812 & 0.908 & -- & REC-WIN & Fine-tuned & VICReg$^{*}$ & SSL \\
Conv + VICReg$^{*}$ + multicorpus   & --    & 0.812 & --    & 0.394 & REC-WIN & Frozen & VICReg$^{*}$ & SSL \\
Conv + multicorpus (FT)             & --    & 0.812 & --    & -- & REC-WIN & Fine-tuned & VICReg$^{*}$ & SSL \\
Conv + Pred-JEPA (20ep)             & --    & 0.808 & 0.899 & -- & REC-WIN & Frozen & Pred-JEPA & SSL \\
Conv + NaiveJEPA (20ep)             & --    & 0.779 & 0.863 & 0.389 & REC-WIN & Frozen & NaiveJEPA & SSL \\
EMA-JEPA on Conv (150ep)            & --    & 0.764 & --    & -- & REC-WIN & Frozen & Pred-JEPA & SSL \\
\midrule
\multicolumn{9}{l}{\textit{LaBraM-style patch Transformer (1.2M parameters)}} \\
LaBraM-style + VICReg$^{*}$         & 0.775$^{a}$ & -- & -- & -- & WIN & Frozen & VICReg$^{*}$ & SSL \\
LaBraM-style + SIGReg               & 0.725$^{a}$ & -- & -- & -- & WIN & Frozen & SIGReg & SSL \\
LaBraM-style + Pred-JEPA (20ep)     & -- & 0.647 & -- & -- & REC-WIN & Frozen & Pred-JEPA & SSL \\
\midrule
\multicolumn{9}{l}{\textit{BIOT-style FFT Transformer (1.2M parameters)}} \\
BIOT-style + VICReg$^{*}$           & 0.775$^{a}$ & -- & -- & -- & WIN & Frozen & VICReg$^{*}$ & SSL \\
BIOT-style + SIGReg                 & 0.783$^{a}$ & -- & -- & -- & WIN & Frozen & SIGReg & SSL \\
BIOT-style + Pred-JEPA (20ep)       & -- & 0.790 & 0.866 & -- & REC-WIN & Frozen & Pred-JEPA & SSL \\
\midrule
\multicolumn{9}{l}{\textit{EEGPT-style channel-pool Transformer (0.8M parameters)}} \\
EEGPT-style + VICReg$^{*}$          & collapsed & -- & -- & -- & WIN & Frozen & VICReg$^{*}$ & SSL \\
EEGPT-style + SIGReg                & collapsed & -- & -- & -- & WIN & Frozen & SIGReg & SSL \\
EEGPT-style + Pred-JEPA (20ep)      & -- & 0.613 & 0.642 & -- & REC-WIN & Frozen & Pred-JEPA & SSL \\
EEGPT-style + VICReg (fixed)        & -- & -- & -- & 0.182 & -- & Frozen & VICReg & SSL \\
EEGPT-style + VICReg (no chan drop) & -- & -- & -- & 0.358 & -- & Frozen & VICReg & SSL \\
\bottomrule
\multicolumn{9}{l}{%
  ``Data'' column: SSL = TUAB SSL training split (2717 recordings, no labels).
  Multicorpus = TUAB+TUEV+TUSZ+TUEP (4$\times$ data).
  BACC\textsubscript{TUEV} = 6-class macro BACC on TUEV event detection (frozen logistic probe).
  Pred-JEPA = GRU predictive objective (collapsed; eval fixed, see \S\ref{sec:bugs}).
  NaiveJEPA = I-JEPA-style temporal split, MSE only, no anti-collapse term.}
\end{tabular}}
\end{table}

\paragraph{Key observations.}
\begin{enumerate}[noitemsep]
  \item \textbf{Conv1D dominates} all encoder families despite fewest parameters (0.4M).
  \item \textbf{Corruption augmentation} (time masking + spike injection) is the
    single most impactful addition: $+0.014$ WIN BACC, $+0.029$ REC BACC.
  \item \textbf{Spectral regularisation at coefficient 0.1} gives the best
    per-recording frozen result (0.836) but was trained with the VICReg bug;
    result must be interpreted with caution.
  \item \textbf{LaBraM-style and BIOT-style} reach the same accuracy (0.775)
    despite fundamentally different tokenisation (raw patches vs.\ FFT magnitude),
    suggesting the bottleneck is encoder capacity at 1.2M params, not tokenisation.
  \item \textbf{EEGPT-style collapsed} under both loss functions with the buggy
    coefficients; rehabilitation with fixed VICReg was pending at write-up time.
  \item \textbf{SIGReg $\geq$ VICReg} for Conv (+0.005 WIN), but
    \textbf{SIGReg $<$ VICReg} for LaBraM-style ($-0.050$ accuracy) ---
    the loss--architecture interaction is non-trivial.
  \item \textbf{Predictive EMA-JEPA} (GRU predictor, frame-prediction objective)
    achieved only 0.764 REC BACC --- the \emph{worst} result of all variants.
    The frame-prediction objective encourages temporal dynamics representations
    rather than class-discriminative statistics.
\end{enumerate}

%% ============================================================
\section{Per-Model Mathematical Descriptions}
\label{sec:math}
%% ============================================================

For each encoder, we provide:
(i) the complete energy function $E$ with explicit loss component definitions
and weight choices; and
(ii) the encoder architecture with its primary design choices.

\subsection{Conv1D Temporal Encoder}
\label{sec:math_conv}

\paragraph{Energy function.}
\begin{equation}
  \boxed{E_{\mathrm{Conv}} = w_2 \cdot L_2 + w_V \cdot L_V + w_C \cdot L_C
    + w_{\mathrm{spec}} \cdot L_{\mathrm{spec}}}
\end{equation}
with weights $(w_2, w_V, w_C) = (25, 25, 1)$ for VICReg and
$w_{\mathrm{spec}} \in \{0, 0.05, 0.1, 0.3, 1.0\}$ for the spectral ablation
(0 = no spectral loss).

\textbf{Component $L_2$ (invariance / reconstruction energy):}
\begin{equation}
  L_2 = \frac{1}{N} \sum_{i=1}^{N} \|\mathbf{z}_{a,i} - \mathbf{z}_{b,i}\|^2,
\end{equation}
where $\mathbf{z}_{a,i}, \mathbf{z}_{b,i} \in \mathbb{R}^{d_p}$ are the
projected embeddings of the two views for sample $i$ in the batch of size $N$.
This is the core energy: minimising $L_2$ pushes both views to the same point.

\textbf{Component $L_V$ (variance / anti-collapse hinge):}
\begin{equation}
  L_V = \frac{1}{2 d_p} \sum_{j=1}^{d_p} \Bigl[
    \max\!\bigl(0,\; \gamma - \mathrm{Std}(\mathbf{Z}_{a, :, j})\bigr) +
    \max\!\bigl(0,\; \gamma - \mathrm{Std}(\mathbf{Z}_{b, :, j})\bigr)
  \Bigr], \quad \gamma = 1.
\end{equation}
$\mathbf{Z}_{a} \in \mathbb{R}^{N \times d_p}$ is the matrix of all projected embeddings
for view $a$; $\mathbf{Z}_{a,:,j}$ is its $j$-th column.
$L_V$ penalises any dimension $j$ whose standard deviation falls below 1,
preventing collapse to a constant vector.

\textbf{Component $L_C$ (covariance / decorrelation):}
\begin{equation}
  L_C = \frac{1}{2 d_p(d_p{-}1)} \Bigl[
    \sum_{i \neq j} \mathrm{Cov}(\mathbf{Z}_a)_{ij}^2 +
    \sum_{i \neq j} \mathrm{Cov}(\mathbf{Z}_b)_{ij}^2
  \Bigr],
\end{equation}
where $\mathrm{Cov}(\mathbf{Z}) = (\mathbf{Z} - \bar{\mathbf{Z}})^\top
(\mathbf{Z} - \bar{\mathbf{Z}}) / (N-1) \in \mathbb{R}^{d_p \times d_p}$.
$L_C$ penalises correlated dimensions, encouraging the encoder to spread
information across independent dimensions.

\textbf{Component $L_{\mathrm{spec}}$ (spectral auxiliary):}
\begin{equation}
  L_{\mathrm{spec}} = \frac{1}{N f} \sum_{i=1}^{N} \sum_{m=1}^{f}
    \Bigl(\log \hat{S}_{i,m} - \log S_{i,m}\Bigr)^2,
\end{equation}
where $S_{i,m}$ is the target Welch power spectral density at frequency bin $m$
for sample $i$, and $\hat{S}_{i,m}$ is predicted by an auxiliary linear head
from the encoder output $\mathbf{z}_i \in \mathbb{R}^{d}$.
This DDSP-inspired loss encourages the encoder to capture spectral structure.

\paragraph{Architecture.}
The Conv1D encoder (from scratch, this work) processes raw EEG
$\mathbf{X} \in \mathbb{R}^{C \times T}$, $C{=}19$, $T{=}2000$.

\begin{equation}
  \mathbf{H}^{(l)} = \mathrm{GELU}\!\Bigl(\mathrm{BN}\!\bigl(
    \mathrm{Conv1d}_{F_l,\, k=7,\, s=2}(\mathbf{H}^{(l-1)})\bigr)\Bigr),
  \quad l = 1,\ldots,4,
\end{equation}
with $F_0 = 19$ (input channels), $F_1 = 64$, $F_2 = 128$, $F_3 = 256$,
$F_4 = 256$.  Each stride-2 convolution halves the time dimension:
$T \to 1000 \to 500 \to 250 \to 125$.  Global average pooling over the time
dimension, followed by a linear projection $256 \to 256$:
\begin{equation}
  \mathbf{z} = \mathbf{W}\! \left(\frac{1}{125} \sum_{t=1}^{125}
    \mathbf{H}^{(4)}_{:,t}\right) \in \mathbb{R}^{256}.
\end{equation}
Total parameters: $\approx 0.4$M.
Projector MLP: $256 \to 1024 \to 1024$ (non-linear, batch-normed), output
$\mathbf{z} \mapsto \hat{\mathbf{z}} \in \mathbb{R}^{1024}$.

\subsection{LaBraM-Style Patch Transformer Encoder}
\label{sec:math_labram}

\paragraph{Energy function.}
Same as Conv1D: $E_{\mathrm{LaBraM}} = 25 L_2 + 25 L_V + L_C$
(or SIGReg: $E = L_2 + \lambda_{\mathrm{SIG}} L_{\mathrm{SIG}}$, see
Section~\ref{sec:math_ssl}).

\paragraph{Architecture.}
Inspired by LaBraM~\cite{labram} without the VQ-codebook neural tokenizer
(from scratch in this work).
Each EEG channel is split into $P{=}10$ non-overlapping patches of
$w{=}200$ timepoints (1 second each).
\begin{equation}
  \mathbf{p}_{c,k} = \mathbf{X}_{c,\, (k-1)w:kw} \in \mathbb{R}^{w},
  \quad c \in [C],\; k \in [P].
\end{equation}
Linear patch embedding: $\mathbf{e}_{c,k} = \mathbf{W}_p \mathbf{p}_{c,k}
+ \mathbf{b}_p \in \mathbb{R}^{d_e}$, $d_e{=}128$.
Learnable temporal embedding $\mathbf{te}_k \in \mathbb{R}^{d_e}$ and
channel (spatial) embedding $\mathbf{se}_c \in \mathbb{R}^{d_e}$:
\begin{equation}
  \hat{\mathbf{e}}_{c,k} = \mathbf{e}_{c,k} + \mathbf{te}_k + \mathbf{se}_c.
\end{equation}
Sequence of $C \times P = 190$ tokens fed to 4 Transformer
blocks~\cite{transformer} with pre-layer normalisation:
\begin{equation}
  \mathrm{Attn}(Q,K,V) = \mathrm{softmax}\!\left(
    \frac{\mathrm{LN}(Q)\,\mathrm{LN}(K)^\top}{\sqrt{d_h}}\right) V,
  \quad d_h = d_e/n_h = 32,\; n_h{=}4.
\end{equation}
FFN hidden dimension: $4 d_e = 512$.
Mean-pool over all 190 tokens: $\bar{\mathbf{e}} = \frac{1}{190}\sum \hat{\mathbf{e}}$.
Linear projection $d_e \to 256$: $\mathbf{z} = \mathbf{W}_\mathrm{out}\bar{\mathbf{e}}$.
Total parameters: $\approx 1.2$M.

\subsection{BIOT-Style FFT Transformer Encoder}
\label{sec:math_biot}

\paragraph{Energy function.}
Same as Conv1D (VICReg or SIGReg).

\paragraph{Architecture.}
Inspired by BIOT~\cite{biot} (from scratch in this work).
The key difference from LaBraM-style is the tokenisation: instead of raw time
samples, each patch $\mathbf{p}_{c,k}$ is embedded via its Fourier magnitude spectrum.
\begin{equation}
  \mathbf{s}_{c,k}[m] = \left|\,\sum_{n=0}^{w-1}
    \mathbf{p}_{c,k}[n]\, e^{-j 2\pi m n / w}\right|,
  \quad m = 0,\ldots,\tfrac{w}{2},
  \quad w{=}200.
\end{equation}
This yields a 101-dimensional frequency-domain feature vector per patch
(the one-sided magnitude spectrum).  Linear projection:
$\mathbf{e}_{c,k} = \mathbf{W}_s \mathbf{s}_{c,k} \in \mathbb{R}^{d_e}$, $d_e{=}128$.
The remaining architecture (channel embedding, position embedding, Transformer)
is identical to LaBraM-style.  The FFT discards phase information, which is why
BIOT-style underperforms on tasks where temporal dynamics (e.g., spike-wave
complexes, burst suppression) are diagnostically relevant.
Total parameters: $\approx 1.2$M.

\subsection{EEGPT-Style Hierarchical Channel-Pool Transformer}
\label{sec:math_eegpt}

\paragraph{Energy function.}
Same as Conv1D (VICReg or SIGReg).

\paragraph{Architecture.}
Inspired by EEGPT~\cite{eegpt} (from scratch in this work).
The key architectural feature is a \emph{spatial channel-pooling step} before the
temporal Transformer.

Each patch $\mathbf{p}_{c,k} \in \mathbb{R}^{w}$ ($w{=}200$) is linearly embedded
with a learnable channel codebook $\mathbf{s}_c$:
\begin{equation}
  \mathrm{tok}_{c,k} = \mathbf{W}_p \mathbf{p}_{c,k} + \mathbf{s}_c \in \mathbb{R}^{d_e},
  \quad d_e{=}128.
\end{equation}

\textit{Spatial channel-pooling (cross-attention).}
A learnable query $\mathbf{q} \in \mathbb{R}^{1 \times d_e}$ attends over all
$C{=}19$ channel tokens at time patch $k$:
\begin{equation}
  \mathbf{g}_k = \mathrm{Attn}\!\left(\mathbf{q};\;
    [\mathrm{tok}_{1,k}, \ldots, \mathrm{tok}_{C,k}]\right)
  \in \mathbb{R}^{d_e}.
\end{equation}
This reduces the $C \times P = 190$ token sequence to $P{=}10$ summary tokens.

\textit{Temporal Transformer.}
4 standard Transformer blocks with learned positional embeddings over the 10 summary
tokens: $\mathbf{h} = \mathrm{TransformerEncoder}([\mathbf{g}_1,\ldots,\mathbf{g}_P])
\in \mathbb{R}^{P \times d_e}$.
Mean-pool and linear projection $d_e \to 256$: $\mathbf{z} \in \mathbb{R}^{256}$.
Total parameters: $\approx 0.8$M.

\paragraph{Failure mode: channel-drop vs.\ channel-pool conflict.}
The channel-pool cross-attention maps 19 channel tokens to 1 summary token per
time patch.  When the \emph{channel-drop augmentation} ($p{=}0.2$ per channel)
independently zeros different channels in view $a$ and view $b$, the resulting
summary tokens $\mathbf{g}_k^a$ and $\mathbf{g}_k^b$ look structurally different
even though they derive from the same underlying signal.  The invariance loss
$L_2 = \|\mathbf{z}_a - \mathbf{z}_b\|^2$ cannot be minimised without the encoder
learning channel-invariant features, which defeats the purpose of the channel
embedding.  This created the EEGPT collapse observed in our experiments.

\subsection{SSL Objectives: VICReg and SIGReg}
\label{sec:math_ssl}

\paragraph{VICReg~\cite{vicreg}.}
Let $\mathbf{Z}, \mathbf{Z}' \in \mathbb{R}^{N \times d_p}$ be the batch of
projected embeddings for views $a$ and $b$.
\begin{equation}
  \boxed{E_{\mathrm{VICReg}} = w_2 \cdot L_2 + w_V \cdot L_V + w_C \cdot L_C}
  \quad\text{with}\; (w_2, w_V, w_C) = (25, 25, 1).
\end{equation}
$L_2$, $L_V$, $L_C$ are defined in Section~\ref{sec:math_conv}.
The paper-correct coefficients $(25, 25, 1)$ from Bardes et al.~\cite{vicreg}
were used in all fixed runs; the buggy runs used $(1, 1, 1)$ (see
Section~\ref{sec:conclusion}).

\paragraph{SIGReg (Batched Characteristic Slicing / BCS).}
SIGReg replaces $L_V$ and $L_C$ with a Gaussianisation surrogate based on the
Epps--Pulley test statistic~\cite{epps1983}.  For $K$ random unit-sphere
projection directions $\mathbf{v}_k \sim \mathcal{U}(\mathcal{S}^{d_p-1})$:
\begin{equation}
  L_{\mathrm{SIG}} = \frac{1}{2K} \sum_{k=1}^{K}
    \Bigl[\mathrm{EP}_N\!\bigl(\{\mathbf{v}_k^\top \mathbf{z}_{a,i}\}_{i}\bigr)
         + \mathrm{EP}_N\!\bigl(\{\mathbf{v}_k^\top \mathbf{z}_{b,i}\}_{i}\bigr)\Bigr],
\end{equation}
where $\mathrm{EP}_N(\{u_i\})$ is the Epps--Pulley test statistic measuring
deviation of the empirical distribution from Gaussian:
\begin{equation}
  \mathrm{EP}_N(\{u_i\}) = \int_{-3}^{3}
    e^{-t^2/2} \left|\hat\varphi_N(t) - e^{-t^2/2}\right|^2 dt,
\end{equation}
where $\hat\varphi_N(t) = \frac{1}{N}\sum_{i=1}^{N} e^{jtu_i}$ is the empirical
characteristic function.  In practice, the integral is approximated by 10
equally-spaced points in $[-3, 3]$.
The full energy is:
\begin{equation}
  \boxed{E_{\mathrm{SIGReg}} = L_2 + \lambda \cdot L_{\mathrm{SIG}}}
  \quad\text{with}\; \lambda = 10,\; K{=}256.
\end{equation}
SIGReg is theoretically stronger than VICReg's variance hinge: representations
collapse iff their random projections become non-Gaussian, which is a more
comprehensive collapse test.  In practice, SIGReg improves Conv1D by 0.005
WIN BACC but \emph{hurts} LaBraM-style ($-0.050$ accuracy), showing that the
loss--architecture coupling is non-trivial.

\subsection{MLP Classifier Architecture}
\label{sec:math_mlp}

The downstream classification probe is a 3-layer MLP:
\begin{align}
  \mathbf{h}_1 &= \mathrm{ReLU}(\mathbf{W}_1 \mathbf{z} + \mathbf{b}_1),
    \quad \mathbf{W}_1 \in \mathbb{R}^{128 \times 256}, \\
  \mathbf{h}_2 &= \mathrm{ReLU}(\mathbf{W}_2 \mathbf{h}_1 + \mathbf{b}_2),
    \quad \mathbf{W}_2 \in \mathbb{R}^{64 \times 128}, \\
  \hat{\mathbf{y}} &= \mathrm{Softmax}(\mathbf{W}_3 \mathbf{h}_2 + \mathbf{b}_3),
    \quad \mathbf{W}_3 \in \mathbb{R}^{2 \times 64}.
\end{align}
Implemented via \texttt{sklearn.neural\_network.MLPClassifier} with
\texttt{hidden\_layer\_sizes=(128, 64)}, early stopping, and
$\ell_2$ regularisation.

In \textbf{frozen mode}, only $\{\mathbf{W}_l, \mathbf{b}_l\}_{l=1}^{3}$
are updated; $f_\theta$ parameters receive no gradient.
In \textbf{fine-tuned mode}, all parameters (encoder + probe) are optimised
jointly with a lower learning rate ($10^{-5}$) for the encoder.

Several probe architectures were considered but not implemented:
logistic regression (sklearn LogReg), which is the standard SSL community probe
and would allow direct comparison to SimCLR/MoCo literature, was never added.
Our MLP is more generous than LogReg, which inflates absolute numbers relative
to published comparisons.

\subsection{Training and Testing Protocol}
\label{sec:protocol}

\paragraph{SSL pre-training (ssl split).}
\begin{itemize}[noitemsep]
  \item Dataset: TUAB SSL training split (2\,717 recordings, labels discarded).
  \item Window extraction: non-overlapping 10-second windows, 2\,000 timepoints
    at 200\,Hz; per-channel z-score normalisation.
  \item Optimiser: AdamW, $\beta_1{=}0.9$, $\beta_2{=}0.999$,
    weight decay $10^{-4}$, learning rate $10^{-4}$ with cosine annealing.
  \item Batch size: 64 windows (128 views per step after two-view augmentation).
  \item Duration: 20 epochs over all training windows.
  \item Augmentation stack (applied independently to each view):
    Gaussian noise ($\sigma{=}0.1$), amplitude jitter ($\pm 20\%$),
    channel dropout (20\% of channels zeroed independently per view),
    time masking (20\% of timepoints zeroed, corruption variant),
    spike injection ($\pm 6\sigma$, 20\% of windows, corruption variant).
\end{itemize}

\paragraph{Probe training (ssl training split with labels).}
\begin{itemize}[noitemsep]
  \item Dataset: 2\,717 TUAB training recordings with normal/abnormal labels.
  \item Feature extraction: $N{=}16$ evenly-spaced 10-second windows per recording
    $\to$ frozen encoder $\to$ 16 embeddings of $\mathbb{R}^{256}$
    $\to$ mean-pool $\to$ one 256-dim vector per recording.
  \item Probe: sklearn MLPClassifier, \texttt{hidden\_layer\_sizes=(128, 64)},
    \texttt{max\_iter=500}, \texttt{early\_stopping=True}.
  \item Standardisation: StandardScaler fitted on training embeddings only;
    same scaler applied to eval embeddings.
\end{itemize}

\paragraph{Evaluation (eval split).}
\begin{itemize}[noitemsep]
  \item Dataset: 276 TUAB evaluation recordings (patient-disjoint from training).
  \item \textbf{Per-window (WIN):} extract all non-overlapping 10-second windows,
    classify each independently, compute BACC over all window predictions.
    Primary metric for benchmarking comparison.
  \item \textbf{Per-recording (REC-WIN):} same $N{=}16$ evenly-spaced window
    extraction as probe training, mean-pool, classify pooled vector.
    Primary metric for clinical deployment evaluation.
  \item No test-time augmentation. No threshold tuning on the eval split.
\end{itemize}

%% ============================================================
\section{Combined Best Results (Per-Recording Focus)}
\label{sec:best}
%% ============================================================

Table~\ref{tab:best} consolidates the best results from the literature and
from this work, filtered to include only per-recording metrics --- the clinically
meaningful unit of analysis for TUAB (one decision per patient).
Per-window results from published SSL models are included for benchmarking
context, clearly labelled.

\begin{table}[htbp]
\centering
\caption{%
  Best results on TUAB v3.0.1 evaluation set.
  Primary focus: per-recording (REC-WIN, one prediction per patient).
  Per-window (WIN) results from published SSL models included for benchmarking.
  \textbf{Bold} = best in each category.
  EEGNet and ShallowConvNet are supervised baselines.
  FEMBA and LaBraM (EEG-Bench) use external large-scale pre-training data.
}
\label{tab:best}
\small
\setlength{\tabcolsep}{5pt}
\resizebox{\textwidth}{!}{%
\begin{tabular}{llccccc}
\toprule
\textbf{Method} & \textbf{Encoder} & \textbf{Params} &
  \textbf{BACC\textsubscript{WIN}} &
  \textbf{BACC\textsubscript{REC}} &
  \textbf{AUROC\textsubscript{REC}} &
  \textbf{Mode} \\
\midrule
\multicolumn{7}{l}{\textit{Supervised baselines (TUAB labels only)}} \\
EEGNet~\cite{eegnet}   & depthwise CNN & 9k & 0.796 & 0.824 & 0.913 & FT \\
ShallowConvNet~\cite{shallowconvnet} & shallow CNN & 44k & 0.777 & 0.803 & 0.893 & FT \\
Deep4Net~\cite{shallowconvnet} & deep CNN & -- & 0.816$^b$ & -- & -- & FT \\
\midrule
\multicolumn{7}{l}{\textit{Foundation models with large external pre-training (published)}} \\
FEMBA-Base~\cite{femba} & Bidirect.\ Mamba & 47.7M & -- & 0.810 & 0.889 & FT (TUEG 21kh) \\
FEMBA-Huge~\cite{femba} & Bidirect.\ Mamba & 386M & -- & \textbf{0.818} & \textbf{0.901} & FT (TUEG 21kh) \\
LaBraM-Base~\cite{labram,eeg_bench} & Patch Transformer & 5.8M & 0.814 & 0.838 & -- & FT \\
\midrule
\multicolumn{7}{l}{\textit{SSL models, TUAB-only pre-training (published)}} \\
ContraWR~\cite{contrawr} & CNN & -- & 0.775 & -- & -- & Frozen \\
BENDR~\cite{bendr} & Conv+Transformer & 4M & 0.717 & -- & -- & FT \\
BIOT~\cite{biot} & FFT Transformer & 3.2M & 0.802 & -- & -- & Frozen \\
LaBraM-Base~\cite{labram} & Patch Transformer & 5.8M & \textbf{0.814} & -- & -- & Frozen \\
\midrule
\multicolumn{7}{l}{\textit{EB-JEPA (this work) --- SSL frozen (TUAB only)}} \\
VICReg base & Conv1D & 0.4M & 0.756 & 0.796 & 0.888 & SSL Frozen \\
VICReg + corruption & Conv1D & 0.4M & 0.770 & 0.825 & 0.904 & SSL Frozen \\
SIGReg + corruption & Conv1D & 0.4M & 0.775 & 0.825 & 0.904 & SSL Frozen \\
VICReg + spectral 0.1$^{\dagger}$ & Conv1D & 0.4M & 0.765 & 0.836 & 0.887 & SSL Frozen \\
\midrule
\multicolumn{7}{l}{\textit{EB-JEPA (this work) --- fine-tuned}} \\
VICReg + corruption $\to$ FT & Conv1D & 0.4M & -- & \textbf{0.812} & \textbf{0.908} & Fine-tuned \\
\bottomrule
\multicolumn{7}{l}{%
  $^b$Accuracy (not BACC); TUAB near-balanced.
  $^{\dagger}$Trained with buggy VICReg coefficients (inv\_coeff=1 vs.\ correct 25).}
\end{tabular}}
\end{table}

\paragraph{Discussion of Table~\ref{tab:best}.}
\begin{itemize}[noitemsep]
  \item The 0.4M Conv1D encoder matches the 5.8M LaBraM-Base at per-window
        BACC when using SIGReg + corruption augmentation (both 0.775).
  \item Per-recording, the fine-tuned encoder reaches BACC
        $0.812 \pm 0.004$ (3 seeds, \emph{final} epoch), statistically
        matching supervised EEGNet ($0.812 \pm 0.013$); the frozen probe
        reaches $0.819 \pm 0.004$.  SSL pre-training without labels thus
        \emph{reaches} --- but does not exceed --- the supervised baseline
        under patient-level aggregation.  (An earlier single-seed,
        \emph{best-epoch} reading of 0.837 suggested SSL beat EEGNet; that
        selected the best epoch on the evaluation set, which peeks at the
        test data, and is retracted --- see Lesson~4.)
  \item The 3.9 pp gap to LaBraM at per-window remains unexplained; the three
        candidate explanations (tokeniser quality, model scale, pretraining
        duration) are confounded at this stage.
  \item AUROC 0.908 (fine-tuned, 3-seed mean) confirms that the representations are
        high-quality even if the linear separability (frozen probe) is
        not yet optimal.
\end{itemize}

%% ============================================================
% Section 5 intentionally reserved
%% ============================================================

%% ============================================================
\section{Conclusion and Future Work}
\label{sec:conclusion}
%% ============================================================

\subsection{Summary}
We presented EB-JEPA, a two-view energy-based SSL framework for clinical EEG
abnormality detection on TUAB.  The primary contribution is not a new SOTA number
but a \emph{systematic diagnostic account} of what works, what fails, and why.

\paragraph{Positive results.}
A 0.4M Conv1D encoder pre-trained for 20 epochs on 2\,717 unlabelled recordings
reaches 0.775 WIN BACC (frozen), matching the published ContraWR result and
within 3.9 pp of LaBraM-Base.  Corruption augmentation (+time masking + spike
injection) is the single most effective ingredient ($+0.014$ WIN, $+0.029$ REC).
Fine-tuning the pre-trained encoder reaches 0.812 REC BACC (3 seeds, final
epoch), matching EEGNet (0.812) with no labels used during pre-training.

\paragraph{Diagnostic findings.}
\begin{enumerate}[noitemsep]
  \item \textbf{VICReg coefficient bug.}
    The default \texttt{inv\_coeff=1.0} was 25$\times$ below the paper-correct value.
    With $(w_2, w_V, w_C) = (1,1,1)$ instead of $(25,25,1)$, the variance term is
    underweighted; the encoder converges to a partially collapsed minimum.
    Diagnostic signature: EEGPT inv\_loss 0.103 (buggy, epoch 19) vs.\ 0.018
    (fixed, epoch 11).

  \item \textbf{Per-window / per-recording evaluation gap.}
    The framework defaults to REC-WIN ($N{=}16$), which is $\approx 5$\,pp above WIN
    for the same encoder.  Comparing REC to WIN across papers produces a spurious
    advantage.  This caused our initial claim of ``beating LaBraM'' (0.836 REC vs.\
    0.814 WIN) to be retracted upon discovering the protocol difference.

  \item \textbf{EEGPT architectural conflict.}
    Channel-pool (19 channels $\to$ 1 token per time patch) combined with
    channel-drop augmentation creates cross-view structural incompatibility:
    the invariance loss $L_2$ cannot be minimised without learning features
    that are independent of which channels are present, which degrades the
    representation for fixed-channel TUAB inference.
\end{enumerate}

\subsection{Open Scientific Questions}
\begin{description}[noitemsep]
  \item[D1] Is EEGPT collapse due to the architectural conflict, the VICReg bug,
    or both independently?  Jobs with fix-only (no channel-drop change) and
    fix + no channel-drop are needed to separate the two.
  \item[D2] What explains the 3.9\,pp gap to LaBraM-Base?
    Candidates: (i) tokeniser quality (LaBraM uses a VQ-codebook pretrained on
    Fourier spectra), (ii) model scale (5.8M vs.\ 0.4M), (iii) pre-training
    duration (200 vs.\ 20 epochs), (iv) data volume.  Each requires a separate
    ablation.
  \item[D3] LogReg vs.\ MLP probe: our MLP gives $\approx$0.01--0.02 BACC higher
    than LogReg would; switching is needed for direct comparison to SimCLR / BYOL
    literature.
  \item[D4] Does spectral regularisation add independent signal at the correct
    VICReg coefficients?  The current best (0.836 REC) used buggy coefficients;
    the combination is untested.
  \item[D5] Does the aperiodic spectral audit~\cite{spectral_audit} apply to
    EB-JEPA?  The flattening test should be run on our Conv1D encoder to determine
    whether SIGReg representations capture oscillatory biomarkers or merely the
    1/$f$ broadband slope.
\end{description}

\subsection{Future Work}

\begin{enumerate}[noitemsep]
  \item \textbf{Neural tokeniser for LaBraM-style encoder.}
    LaBraM's pretrained VQ-codebook neural spectrum tokeniser (trained via masked
    patch prediction on Fourier spectra~\cite{labram}) is the most likely
    differentiator at 5.8M params.  Training this tokeniser separately and plugging
    into our 1.2M backbone is the highest-priority experiment for closing the gap.

  \item \textbf{Longer pre-training and larger scale.}
    All runs used 20 epochs.  LaBraM used $\approx$200 epochs on 2\,500h of diverse
    EEG.  FEMBA~\cite{femba} used 21,000h.  Doubling to 40 epochs on TUAB-only
    is a cheap test of whether capacity or duration limits current performance.

  \item \textbf{EEGPT with encoder-specific augmentation policy.}
    Removing channel-drop from EEGPT's augmentation stack (while keeping it for
    Conv1D and LaBraM) is the minimal fix to resolve the architectural conflict.

  \item \textbf{SSM-based encoder (Mamba-style).}
    FEMBA~\cite{femba} demonstrates that bidirectional Mamba blocks scale linearly
    in sequence length (vs.\ $\mathcal{O}(N^2)$ for Transformers), which is
    attractive for long clinical EEG recordings.  Replacing the Conv1D backbone
    with a small Mamba encoder ($\approx$7.8M, FEMBA-Tiny scale) would test
    whether the SSM architectural prior improves the energy-based objective.

  \item \textbf{Aperiodic reliance audit.}
    Apply Bindra \& Panwar's phase-preserving Fourier intervention
    framework~\cite{spectral_audit} to all our frozen encoders to determine
    whether EB-JEPA representations exploit the 1/$f$ aperiodic shortcut.

  \item \textbf{Predictive JEPA with per-frame probe.}
    The EMA-GRU predictive variant (BACC 0.764 REC) was probed via global
    average pooling, which discards the temporal structure the objective learned.
    A per-frame temporal probe (one label per time step, aggregated) may recover
    useful information from predictive representations.
\end{enumerate}

\subsection{Methodological Lessons}

This project illustrates four recurring failure modes in EEG deep learning:
\begin{enumerate}[noitemsep]
  \item \textbf{Comparison without protocol matching.}
    WIN vs.\ REC-WIN is the TUAB-specific instance of a general problem:
    different papers measure different quantities.  Always verify the exact
    evaluation protocol of the paper you compare against before reporting a gap.

  \item \textbf{Loss coefficient bugs are invisible from convergence curves.}
    A 25$\times$ underweighted term may still converge to a ``good enough''
    minimum; the bug is only detectable at downstream evaluation or by auditing
    the raw loss magnitudes against the paper's reported ranges.

  \item \textbf{Augmentation--architecture coupling.}
    Augmentation choices that are benign for one architecture can be
    catastrophic for another (channel-drop + channel-pool).
    The augmentation stack should be designed or ablated per encoder family,
    not applied uniformly.
  \item \textbf{Best-epoch selection peeks at the test set.}
    Reporting the epoch with the highest \emph{evaluation} accuracy uses the
    test set for model selection --- a subtle leak.  Our fine-tuned probe
    appeared to reach 0.837 REC BACC at its best epoch but only
    $0.812 \pm 0.004$ at the final epoch (3 seeds); the 2.5\,pp gap is pure
    peeking, and it was exactly the margin that made SSL look like it
    \emph{beat} EEGNet rather than tying it.  Every model must share one fixed
    stopping rule (we report final-epoch, seed-averaged); selecting per-model
    best epochs against the same split you report on manufactures wins that
    vanish under a held-out protocol.
\end{enumerate}

\bigskip
\noindent\textbf{Code and data.}
All encoder architectures, training scripts, evaluation protocols, and the
science-check protocol are available in the project repository
(\texttt{Hack\_the\_worlds/}).
The repository includes the EB-JEPA library~\cite{ebjepa_lib} and all hackathon
experiment logs (\texttt{reports/EXPERIMENTS.md},
\texttt{reports/EVALUATION\_ANALYSIS.md}).

\bibliographystyle{plainnat}
\bibliography{references}

\end{document}
